\section{Related work}
\label{sec:related}

\paragraph{Type encodings for polymorphic logics.}
The systematic study of sound type erasure begins with Meng and
Paulson~\cite{MengPaulson2008}, who documented both the performance cost of
full type annotations and the concrete unsoundness of dropping them
(\cref{prop:tight-uncovered,prop:tight-nil} are fragment forms of their
examples), and proposed partial encodings validated empirically.  Couchot
and Lescuyer~\cite{CouchotLescuyer2007} and Bobot and
Paskevich~\cite{BobotPaskevich2011} gave provably sound guard- and
tag-based encodings of polymorphic FOL into untyped and many-sorted FOL.
Claessen, Lillieström and
Smallbone~\cite{ClaessenLilliestromSmallbone2011} introduced
monotonicity-based erasure for many-sorted logic; their ``calculus 2'', in
which an arbitrary predicate may serve as the guard of a sort provided it
can consistently be extended to be false outside it, is the direct formal
ancestor of our covers.  Our \cref{def:MN} builds the false extension into
the intended model once and for all --- an ill-typed atom is false because
an ill-typed application is not a provable proposition --- so no
per-problem satisfiability check is needed; the price is that atoms which
are not false-extended in $\MN$ (equations, non-rigid $\T$-atoms) can
never serve as covers, which \cref{sec:tightness} shows is essential
anyway.  Blanchette, B\"ohme, Popescu and
Smallbone~\cite{BlanchetteBohmePopescuSmallbone2016} developed the most
comprehensive framework: their \emph{covers} and \emph{inferable type
arguments} correspond to our \cref{def:erasable,lem:inference-unique},
their \emph{featherweight guards} $\mathsf{g}{?}{?}$ --- guards asserted
only for naked variables --- are the satisfiability-preserving relative of
our per-formula equivalence, and their monotonicity calculi point to a
third, complementary omission principle that our development does not
cover (see below).  Their encodings, implemented in Sledgehammer and, in
guard form, in Why3~\cite{FilliatrePaskevich2013}, are global problem
transformations proved sound by model constructions similar in spirit to
\cref{subsec:intended}; the difference in our setting is dependency ---
the forced type $\enc{\tau^p_i}[\tup s/\tup\alpha]$ of an argument
position may mention the other arguments --- and the conclusion is
correspondingly finer-grained: a per-occurrence logical equivalence
relative to $\Comp$ rather than the equisatisfiability of a whole-problem
encoding.

\paragraph{Hammers and CIC.}
For the hammer context see the survey~\cite{BlanchetteKaliszykPaulsonUrban2016}
and the systems for HOL~\cite{PaulsonBlanchette2010,KaliszykUrban2014} and
Mizar~\cite{KaliszykUrban2015}.  CoqHammer~\cite{CzajkaKaliszyk2018}
defined the translation whose optimization we study; its practically
tuned variant already omits selected guards and justifies the choices
empirically, by counterexample-driven testing and reconstruction rates.
The theoretical companion~\cite{Czajka2018} proves such a translation of
a class of pure type systems sound \emph{proof-theoretically} --- any FOL
proof from the translated axioms maps to a source derivation --- but for
the guarded variant only; omission of redundant guards is stated there as
a conjecture, and the proof hinges on an invariant that guards are the
sole channel through which typability information flows into a FOL proof.
Our results give the classical, model-theoretic counterpart for a
fragment with inductive types, and suggest a route to the proof-theoretic
version: \cref{cor:no-loss} shows that a proof from guard-omitted axioms
converts, purely in FOL, into a proof from the guarded axioms plus
$\Comp$, so it suffices to extend the invariant to completion axioms ---
whose source reading is exactly typing inversion: from a reconstruction
of a proof of $p\,\tup\sigma\,\tup t$ one extracts typing evidence for
$\tup t$.  Carrying this out in the setting of~\cite{Czajka2018}
(which would first need extending that framework to inductive types)
remains future work.  Mizar's \emph{conditional cluster registrations},
exported to ATPs by MPTP~\cite{Urban2006}, are Horn sentences of exactly
the completion-axiom shape (``every $\mathit{cardinal}$ is an
$\mathit{ordinal}$''), used at scale in the MizAR
hammer~\cite{KaliszykUrban2015}; our $\Comp$ can be read as the
automatically generated cluster theory of the translated environment.
In the Lean ecosystem, Lean-auto~\cite{LeanAuto2025} faces the same
design choice for its monomorphizing translation.  On the source side,
our erasability criterion parallels Letouzey's program
extraction~\cite{Letouzey2003}: erasable type arguments are those
recoverable from the extracted (term-level) data, and result-only
arguments are exactly the positions where extraction must insert an
unchecked coercion.

\paragraph{Limitations and future work.}
The fragment \Frag\ isolates the polymorphic first-order core of CIC, and
three omissions matter in practice.  (i) \emph{Conversion}: CoqHammer
emits lambda-lifted definitional equations \emph{unguarded}, justified by
a different intended semantics in which equality is convertibility of
(possibly ill-typed) terms; our models interpret equality as identity on
tagged values and junk trees, so unguarded definitional equations are
outside the present theorems (and indeed our criterion conservatively
keeps their guards unless a cover is present).  Reconciling the two
semantics --- a quotient of the junk universe by an untyped conversion
--- is the natural next step.  (ii) \emph{Dependent data and higher-order
features}: type formers with term indices, function types as data, and
non-prenex polymorphism all break the freeness arguments in
\cref{lem:subst-inj} and \cref{def:MN}; hammer translations handle them
by lifting and collapsing, and the interaction of those devices with
omission (in particular, which lifted atoms are false-extended and hence
covers) is open --- the known constraint being that atoms headed by
lifted proposition constants must \emph{not} count as covers, since their
defining equivalences quantify unguarded.  (iii) \emph{Monotonicity}: for
types whose translated theory cannot distinguish domain sizes, guards can
be omitted with \emph{no} cover at all
\cite{ClaessenLilliestromSmallbone2011,BlanchetteBohmePopescuSmallbone2016};
adapting monotonicity inference to a dependently typed source, where the
finiteness of a type can depend on term parameters, is the open problem
already posed in~\cite{Czajka2018}, and our
\cref{prop:tight-uncovered} (a one-element type) marks exactly the
non-monotonic case that any such analysis must detect.

\section{Conclusion}
\label{sec:conclusion}

We have given syntactic criteria for the two dominant redundancies of
guarded hammer translations of CIC --- type guards covered by a strict
hypothesis atom, and type arguments recoverable from retained term
arguments --- and proved them sound, and in the guard case conservative,
for a polymorphic first-order fragment of CIC with inductive definitions.
The proofs are self-contained model constructions; the single point of
contact with CIC metatheory is the (cited) existence of a classical model
of the source environment, and the single semantic insight doing the work
is that the intended interpretation of a translated predicate is
\emph{false-extended}: an ill-typed application is not a provable
proposition.  The criteria evaluate on the emitted first-order formulas
using only the declared types of the translated constants, and are
therefore implementable as a post-processing pass in an existing hammer;
the completion axioms $\Comp$, themselves sound to emit, make guard
omission lossless.  The tightness examples delimit what any extension
must respect: covers must be false-extended atoms, existential guards are
facts, inference needs a retained witness, and predicates must keep the
arguments that make them strict.
